% Part: proof-theory
% Chapter: cut-elimination
% Section: interpolation

\documentclass[../../../include/open-logic-section]{subfiles}

\begin{document}

\olfileid{pt}{cut}{itp}

\olsection{Lema Maehara lan Teorema Interpolasi Craig}

Teorema interpolasi saka William Craig mratélakaké yèn
$!A \Entails !B$, banjur ana !!a{sentence}~$!C$
(\emph{interpolan}) kang marakaké $!A \Entails !C$ lan
$!C \Entails !B$. Kang wigati, interpolan~$!C$ bisa
dipilih saéngga kabèh !!{constant}s lan !!{predicate}s
ing~$!C$ uga ana ing~$!A$ lan~$!B$.
(Kita nganggep ora ana !!{function}s.) Tanpa watesan iki,
$!A$~lan $!B$ bisa kanthi gampang dadi interpolan.

Teorema interpolasi bener kanggo logika orde kapisan klasik
lan tau diarani ``sipat wigati pungkasan saka logika orde
kapisan'' kang ditemokaké. Teorema iki duwé panggunaan
wigati ing teori modèl lan panalaran otomatis.
Motivasi lan panggunaan wiwitané ana ing filsafat ilmu,
amarga bisa dipigunakaké kanggo nyinaoni istilah primitif
kang bisa utawa ora bisa dadi dhasar aksiomatisasi sawijining
teori. Teorema iki uga ndadèkaké teorema definabilitas Beth,
kang mratélakaké yèn sawijining teori bisa nemtokaké konsèp
kanthi implisit, teori iku uga bisa nemtokaké kanthi eksplisit.

Kaya akèh asil ing logika matématis, teorema interpolasi
bisa dibuktèkaké nganggo cara teori modèl utawa teori bukti.
Kaluwihan cara teori bukti yaiku ora mung nuduhaké anané
interpolan, nanging uga maringi algoritma kanggo ngétung
interpolan saka !!a{proof} yèn $!A \Entails !B$,
tuladhané !!a{proof} ing~\Log{G3c} kanggo
$!A \Sequent !B$.

Upama $\Lang L(!A)$ himpunan !!{constant}s lan
!!{predicate}s ing~$!A$, sarta
$\Lang L(\Gamma) = \bigcup \Setabs{\Lang L(!A)}{!A \in
\Gamma}$. Ing sakabèhé pérangan iki, kita nganggep
!!{formula}s kang ora ngemot !!{function}s.

\begin{thm}[Teorema Interpolasi Craig]\ollabel{thm:interpolation}
Upama !!{language} saka~$!A$ yaiku $\Lang L(!A)$ lan
basa saka $!B$ yaiku $\Lang L(!B)$.
Yèn $\Log{G3c} \Proves !A \Sequent !B$, banjur ana
!!a{formula}~$!C$ ing !!{language}~$\Lang L(!C) \subseteq \Lang
L(!A) \cap \Lang L(!B)$ saéngga
$\Log{G3c} \Proves !A \Sequent !C$ lan
$\Log{G3c} \Proves !C \Sequent !B$.
\end{thm}

Kita mbuktèkaké generalisasi teorema interpolasi supaya bisa
migunakaké struktur !!{proof}s ing~\Log{G3c}. Anggepen
antesèden lan suksèden saben sekèn dipérang dadi rong
pérangan. Kita nulis
$\maeh{\Gamma_1}{\Gamma_2} \Sequent
\maeh{\Delta_1}{\Delta_2}$ kanggo nandhakaké yèn
$\Gamma_1$ lan $\Delta_1$ ana ing pérangan kapisan,
dene $\Gamma_2$ lan $\Delta_2$ ana ing pérangan kapindho.
Teorema interpolasi banjur langsung ndhèrèk saka lema iki.

\begin{lem}[Lema Maehara]\ollabel{lem:maehara}
Yèn $\Log{G3c}
\Proves \maeh{\Gamma_1}{\Gamma_2} \Sequent \maeh{\Delta_1}{\Delta_2}$,
banjur ana !!a{formula}~$!C$ ing !!{language}~$\Lang L(!C)
\subseteq \Lang L(\Gamma_1,
\Delta_1) \cap \Lang L(\Gamma_2, \Delta_2)$ saéngga
$\Log{G3c} \Proves \Gamma_1 \Sequent \Delta_1, !C$ lan
$\Log{G3c} \Proves !C, \Gamma_2 \Sequent \Delta_2$.
\end{lem}

\begin{proof}
  Anggepen $\pi \Proves \maeh{\Gamma_1}{\Gamma_2} \Sequent
  \maeh{\Delta_1}{\Delta_2}$. Pratelan iki kita buktèkaké
  kanthi induksi tumrap $n=\pheight{\pi}$.

  Yèn $n=0$, banjur $\maeh{\Gamma_1}{\Gamma_2} \Sequent
  \maeh{\Delta_1}{\Delta_2}$ iku aksioma, lan ana
  !!{formula}~$!A$ atomik ing $\Gamma_1, \Gamma_2$
  uga ing $\Delta_1, \Delta_2$. Ana patang kahanan,
  manut $!A$~ana ing pérangan kapisan utawa kapindho
  ing sisih kiwa, lan ing pérangan kapisan utawa kapindho
  ing sisih tengen.
  \begin{enumerate}
    \item $!A \in \Gamma_1$ lan $!A \in \Delta_1$.
    Kita nggoleki~$!C$ saéngga
    $\Gamma_1 \Sequent \Delta_1, !C$ lan
    $!C, \Gamma_2 \Sequent \Delta_2$ kaloroné !!{provable}.
    Sekèn kapisan wis aksioma tanpa gumantung pilihan~$!C$,
    amarga $!A$ ana ing kiwa lan tengen. Supaya
    $!C, \Gamma_2 \Sequent \Delta_2$ mesthi !!{provable},
    pilihané yaiku $!C \ident \lfalse$.
    \item $!A \in \Gamma_1$ lan $!A \in \Delta_2$.
    Banjur $\Gamma_1 \Sequent \Delta_1, !A$ lan
    $!A, \Gamma_2 \Sequent \Delta_2$ kaloroné aksioma,
    mula kita bisa milih $!C \ident !A$.
    \item $!A \in \Gamma_2$ lan $!A \in \Delta_1$.
    Banjur $!A, \Gamma_1 \Sequent \Delta_1$ lan
    $\Gamma_2 \Sequent \Delta_2, !A$ dadi aksioma,
    nanging~$!A$ ana ing sisih kang kliru. Mula $!A$
    dudu interpolané. Nanging nganggo $\RightR{\lnot}$
    lan $\LeftR{\lnot}$ kita bisa !!{prove}
    $\Gamma_1 \Sequent \Delta_1, \lnot!A$ lan
    $\lnot !A, \Gamma_2 \Sequent \Delta_2$.
    \item $!A \in \Gamma_2$ lan $!A \in \Delta_2$.
    Latihan.
  \end{enumerate}
  Kahanan nalika $\maeh{\Gamma_1}{\Gamma_2} \Sequent
  \maeh{\Delta_1}{\Delta_2}$ dadi aksioma amarga
  $\lfalse \in \Gamma_1$ utawa $\lfalse \in \Gamma_2$
dipasrahaké dadi latihan.

  Yèn $n>0$, $\pi$ pungkasané nganggo inferènsi logis.
  Kahanané dibedakaké manut aturan kang dienggo lan
  pérangan endi kang ngemot formula utama. Ing saben
  kahanan, kita bisa migunakaké hipotesis induksi yèn
  lema iki bener kanggo kabèh !!{proof}s~$\pi_1$ kanthi
  $\pheight{\pi_1} < n$, senadyan sekèn pungkasan~$\pi_1$
dipérang dadi rong pérangan kanthi cara apa waé.
  Gatekna sawatara kahanan tartamtu.
  \begin{enumerate}
    \item $\pi$ pungkasané nganggo $\RightR{\lnot}$ kanthi
    formula utama~$\lnot !A$. Ana rong kahanan: manut
    $\lnot !A$ ana ing pérangan kapisan utawa kapindho,
    $\pi$ pungkasané nganggo salah siji saka
    \[
    \AxiomC{}
    \RightLabel{$\pi_1$}
    \Deduce$\maeh{!A, \Gamma_1}{\Gamma_2} \fCenter
      \maeh{\Delta_1}{\Delta_2}$
    \RightLabel{\RightR{\lnot}}
    \UnaryInf$\maeh{\Gamma_1}{\Gamma_2} \fCenter
      \maeh{\Delta_1, \lnot !A}{\Delta_2}$
    \DisplayProof
    \]
    utawa
    \[
    \AxiomC{}
    \RightLabel{$\pi_1$}
    \Deduce$\maeh{\Gamma_1}{!A, \Gamma_2} \fCenter
      \maeh{\Delta_1}{\Delta_2}$
    \RightLabel{\RightR{\lnot}}
    \UnaryInf$\maeh{\Gamma_1}{\Gamma_2} \fCenter
      \maeh{\Delta_1}{\Delta_2, \lnot !A}$
    \DisplayProof
    \]
    Gatekna yèn nalika !!{formula} utama~$\lnot !A$
    ana ing pérangan kapisan, !!{formula} aktif~$!A$
    uga dilebokaké ing pérangan kapisan premis. Iki
    pilihan kang bisa kita gawé. Hipotesis induksi uga
    bisa ditrapaké yèn formula aktif mau dilebokaké
    ing pérangan liyané, nanging banjur kita ora bisa
    nemokaké interpolan (latihan: cobanen).

    Gatekna dhisik kahanan kapisan nalika $\lnot !A$~ana
    ing pérangan kapisan. Miturut hipotesis induksi,
    ana interpolan kanggo sekèn pungkasan~$\pi_1$,
    yaiku !!a{formula}~$!C_1$ saéngga kaloroné
    \begin{align*}
    !A, \Gamma_1 & \Sequent \Delta_1, !C_1 &&\text{lan} &
    !C_1, \Gamma_2 & \Sequent \Delta_2
    \intertext{bisa dibuktèkaké. Kang dibutuhaké yaiku !!a{formula}~$!C$ saéngga
    rong sekèn iki !!{provable}:}
    \Gamma_1 & \Sequent \Delta_1, \lnot !A, !C &&\text{lan} &
    !C, \Gamma_2 & \Sequent \Delta_2
    \end{align*}
    Cetha yèn kita bisa njupuk $!C \ident !C_1$:
    hipotesis induksi wis njamin sekèn tengen !!{provable},
    lan sekèn kiwa ndhèrèk saka sekèn ing hipotesis induksi
    kanthi~$\RightR{\lnot}$. Hipotesis induksi uga maringi
    $\Lang L(!C_1) \subseteq \Lang L(!A, \Gamma_1,
    \Delta_1) \cap \Lang L(\Gamma_2, \Delta_2)$.
    Amarga $\Lang L(\lnot !A) = \Lang L(!A)$ lan
    $\Lang L(!C) = \Lang L(!C_1)$, kita ngolèhké
    $\Lang L(!C_1) \subseteq \Lang L(\Gamma_1, \Delta_1, \lnot !A) \cap \Lang
    L(\Gamma_2, \Delta_2)$.

    Ing kahanan kapindho, $\lnot !A$~ana ing pérangan
    kapindho. !!{formula} aktif ing premis uga kita
    lebokaké ing pérangan kapindho, banjur hipotesis
    induksi ditrapaké marang
    \begin{align*}
    \Gamma_1 & \Sequent \Delta_1, !C_1 &&\text{lan} &
    !C_1, !A, \Gamma_2 & \Sequent \Delta_2
    \intertext{kang bisa dibuktèkaké. Yèn aturan \RightR{\lnot} ditrapaké marang sekèn kapindho, kita ngolèhké !!{proof}s kanggo}
    \Gamma_1 & \Sequent \Delta_1, !C_1 &&\text{lan} &
    !C_1, \Gamma_2 & \Sequent \Delta_2, \lnot !A
    \end{align*}
    Iki pancèn sekèn-sekèn kang kudu !!{provable} supaya
    $!C_1$ dadi interpolan; kita njupuk manèh
    $!C \ident !C_1$. Kaya sadurungé,
    $\Lang L(!C_1) \subseteq \Lang L(\Gamma_1, \Delta_1) \cap \Lang L(\Gamma_2,
    \Delta_2, \lnot !A)$.
    \item $\pi$ pungkasané nganggo $\RightR{\lif}$ kanthi
    !!{formula} utama~$!A \lif !B$. Ana manèh rong kahanan,
    manut $!A \lif !B$ ana ing pérangan kapisan utawa
    kapindho. !!{proof}~$\pi$ pungkasané nganggo salah siji saka
    \[
    \AxiomC{}
    \RightLabel{$\pi_1$}
    \Deduce$\maeh{!A, \Gamma_1}{\Gamma_2} \fCenter
      \maeh{\Delta_1, !B}{\Delta_2}$
    \RightLabel{\RightR{\lif}}
    \UnaryInf$\maeh{\Gamma_1}{\Gamma_2} \fCenter
      \maeh{\Delta_1, !A \lif !B}{\Delta_2}$
    \DisplayProof
    \]
    utawa
    \[
    \AxiomC{}
    \RightLabel{$\pi_1$}
    \Deduce$\maeh{\Gamma_1}{!A, \Gamma_2} \fCenter
      \maeh{\Delta_1}{\Delta_2, !B}$
    \RightLabel{\RightR{\lif}}
    \UnaryInf$\maeh{\Gamma_1}{\Gamma_2} \fCenter
      \maeh{\Delta_1}{\Delta_2, !A \lif !B}$
    \DisplayProof
    \]
    Manèh, !!{formula}s aktif dilebokaké ing pérangan
    premis kang padha karo pérangan !!{formula} utama
    ing konklusi. Ing kahanan kapisan, hipotesis induksi
    maringi !!{proof}s kanggo
    \begin{align*}
    !A, \Gamma_1 & \Sequent \Delta_1, !B, !C_1 &&\text{lan} & !C_1, \Gamma_2 & \Sequent \Delta_2
    \intertext{Sekèn kiwa dadi premis kang cocog kanggo \RightR{\lif}. Mula sekèn-sekèn iki !!{provable}:}
    \Gamma_1 & \Sequent \Delta_1, !A \lif !B, !C_1 &&\text{lan} & !C_1, \Gamma_2 & \Sequent \Delta_2
    \end{align*}
    Tegesé, $!C_1$ uga interpolan kanggo sekèn pungkasan
    saka~$\pi$, mula kita bisa njupuk $!C \ident !C_1$ manèh.
    
    Kahanan kapindho ditangani kanthi cara kang padha.
    \item $\pi$ pungkasané nganggo $\LeftR{\lif}$ kanthi
    !!{formula} utama~$!A \lif !B$. Yèn $!A \lif !B$
    ana ing pérangan kapisan, !!{proof}~$\pi$ pungkasané kaya mangkéné:
    \[
    \AxiomC{}
    \RightLabel{$\pi_1$}
    \Deduce$\maeh{\Gamma_1}{\Gamma_2} \fCenter
      \maeh{\Delta_1, !A}{\Delta_2}$
    \AxiomC{}
    \RightLabel{$\pi_2$}
    \Deduce$\maeh{!B, \Gamma_1}{\Gamma_2} \fCenter
      \maeh{\Delta_1}{\Delta_2}$
    \BinaryInf$\maeh{!A \lif !B, \Gamma_1}{\Gamma_2} \fCenter
      \maeh{\Delta_1}{\Delta_2}$
    \DisplayProof
    \]
    Hipotesis induksi saiki maringi patang sekèn
    kang !!{provable}: loro kanggo interpolan $!C_1$
    saka premis kiwa lan loro kanggo interpolan $!C_2$
    saka premis kapindho:
    \begin{align*}
    \Gamma_1 & \Sequent \Delta_1, !A, !C_1 &&\text{(1)} &
    !C_1, \Gamma_2 & \Sequent \Delta_2 && \text{(2)}\\
    !B, \Gamma_1 & \Sequent \Delta_1, !C_2 &&\text{(3)} &
    !C_2, \Gamma_2 & \Sequent \Delta_2 && \text{(4)}
    \end{align*}
    Sekèn (1) lan (3) bisa dadi premis~\LeftR{\lif}
    yèn nganggo pelemahan kanggo nambah $!C_2$ ing sisih
    tengen~(1) lan $!C_1$ ing sisih tengen~(3):
    \[
    \AxiomC{}
    \Deduce$\Gamma_1 \fCenter \Delta_1, !A, !C_1, !C_2$
    \AxiomC{}
    \Deduce$!B, \Gamma_1 \fCenter \Delta_1, !C_1, !C_2$
    \RightLabel{\LeftR{\lif}}
    \BinaryInf$!A \lif !B, \Gamma_1 \fCenter \Delta_1, !C_1, !C_2$
    \DisplayProof
    \]
    Kanthi ngetrapaké $\RightR{\lor}$ kita ngolèhké sekèn
    \[
    !A \lif !B, \Gamma_1 \fCenter \Delta_1, !C_1 \lor !C_2
    \]
    Iki nuduhaké yèn $!C_1 \lor !C_2$ bisa dadi interpolan
    ing kahanan iki. Pancèn, kita bisa !!{prove} sekèn
    liyané kang dibutuhaké saka sekèn (2) lan~(4):
    \[
    \AxiomC{}
    \Deduce$!C_1, \Gamma_2 \fCenter \Delta_2$
    \AxiomC{}
    \Deduce$!C_2, \Gamma_2 \fCenter \Delta_2$
    \RightLabel{\LeftR{\lor}}
    \BinaryInf$!C_1 \lor !C_2, \Gamma_2 \fCenter \Delta_2$
    \DisplayProof
    \]
    Isih perlu dipriksa yèn $!C_1 \lor !C_2$ ana ing
    basa kang trep. Hipotesis induksi maringi
    \begin{align*}
    \Lang L(!C_1) & \subseteq
    \Lang L(\Gamma_1, \Delta_1, !A) \cap \Lang L(\Gamma_2, \Delta_2) &\text{lan}\\
    \Lang L(!C_2) & \subseteq
    \Lang L(!B, \Gamma_1, \Delta_1) \cap \Lang L(\Gamma_2, \Delta_2)
    \intertext{Amarga}
    \Lang L(!A) & \subseteq \Lang L(!A \lif !B) &\text{lan}\\
    \Lang L(!B) & \subseteq \Lang L(!A \lif !B)
    \intertext{kita ngolèhké kaloroné}
    \Lang L(!C_1) & \subseteq
    \Lang L(\Gamma_1, \Delta_1, !A \lif !B) \cap \Lang L(\Gamma_2, \Delta_2) &\text{lan}\\
    \Lang L(!C_2) & \subseteq
    \Lang L(\Gamma_1, \Delta_1, !A \lif !B) \cap \Lang L(\Gamma_2, \Delta_2)
    \intertext{mula, $\Lang L (!C_1 \lor !C_2) = {}$}
    \Lang L(!C_1) \cup \Lang L(!C_2) & \subseteq
    \Lang L(\Gamma_1, \Delta_1, !A \lif !B) \cap \Lang L(\Gamma_2, \Delta_2),
    \end{align*}
    kaya kang dibutuhaké.

    Yèn $!A \lif !B$ ana ing pérangan kapindho,
    kahanané béda nanging ditangani kanthi cara kang
    padha. Hipotesis induksi maringi manèh patang sekèn
    kang !!{provable}:
    \begin{align*}
    \Gamma_1 & \Sequent \Delta_1, !C_1 &&\text{(1)} &
    !C_1, \Gamma_2 & \Sequent \Delta_2, !A && \text{(2)}\\
    \Gamma_1 & \Sequent \Delta_1, !C_2 &&\text{(3)} &
    !C_2, !B, \Gamma_2 & \Sequent \Delta_2 && \text{(4)}
    \end{align*}
    Saiki sekèn (2) lan (4)---kanthi sawatara !!{formula}s
    kang ditambahaké liwat pelemahan---dadi premis~\LeftR{\lif}:
    \[
    \AxiomC{}
    \Deduce$!C_1, !C_2, \Gamma_2 \fCenter \Delta_2, !A$
    \AxiomC{}
    \Deduce$!B, !C_1, !C_2, \Gamma_2 \fCenter \Delta_2$
    \RightLabel{\LeftR{\lif}}
    \BinaryInf$!A \lif !B, !C_1, !C_2, \Gamma_2 \fCenter \Delta_2$
    \DisplayProof
    \]
    Kita perlu ngowahi iki dadi sekèn kang nduwèni siji
    !!{formula} saka $!C_1$ lan $!C_2$, nanging saiki
    ana ing antesèden (amarga iki pérangan kapindho).
    Trapna $\LeftR{\land}$; interpolané bisa dijupuk
    $!C \ident (!C_1 \land !C_2)$. Sekèn (1) lan (3)
    kang isih ana bisa digunakaké kanggo !!{prove}
    sekèn kang cocog kanggo pérangan liyané:
    \[
    \AxiomC{}
    \Deduce$\Gamma_1 \fCenter \Delta_1, !C_1$
    \AxiomC{}
    \Deduce$\Gamma_1 \fCenter \Delta_1, !C_2$
    \RightLabel{\RightR{\land}}
    \BinaryInf$\Gamma_1 \fCenter \Delta_1, !C_1 \land !C_2$
    \DisplayProof
    \]
    Kaya sadurungé, kita nduwèni
    $\Lang L(!C_1 \land !C_2) \subseteq \Lang
    L(\Gamma_1, \Delta_1) \cap \Lang L(\Gamma_2, \Delta_2, !A \lif
    !B)$.

    \item $\pi$ pungkasané nganggo $\RightR{\lforall}$ kanthi
    !!{formula} utama~$\lforall[x][!A(x)]$:
    \[
    \AxiomC{}
    \RightLabel{$\pi_1$}
    \Deduce$\maeh{\Gamma_1}{\Gamma_2} \fCenter
    \maeh{\Delta_1, !A(c)}{\Delta_2}$
    \RightLabel{\RightR{\lforall}}
    \UnaryInf$\maeh{\Gamma_1}{\Gamma_2} \fCenter
      \maeh{\Delta_1, \lforall[x][!A(x)]}{\Delta_2}$
    \DisplayProof
    \]
    utawa
    \[
    \AxiomC{}
    \RightLabel{$\pi_1$}
    \Deduce$\maeh{\Gamma_1}{\Gamma_2} \fCenter
      \maeh{\Delta_1}{\Delta_2, !A(c)}$
    \RightLabel{\RightR{\lforall}}
    \UnaryInf$\maeh{\Gamma_1}{\Gamma_2} \fCenter
      \maeh{\Delta_1}{\Delta_2, \lforall[x][!A(x)]}$
    \DisplayProof
    \]
    Ing kahanan kapisan, hipotesis induksi maringi !!{proof}s
    kanggo $\Gamma_1 \Sequent \Delta_1, !A(c), !C_1$ lan
    $!C_1, \Gamma_2 \Sequent \Delta_2$.
    Kita bisa ngetrapaké $\RightR{\lforall}$ marang kang
    kapisan kanggo ngolèhké
    $\Gamma_1 \Sequent \Delta_1, \lforall[x][!A(x)], !C_1$.
    (Syarat variabel eigen kacukupan, amarga wis kacukupan
    ing aturan \RightR{\lforall} ing pungkasan~$\pi$.)
    Mula $!C_1$ dadi interpolan yèn
    $\Lang L(!C_1) \subseteq \Lang L(\Gamma_1, \Delta_1, \lforall[x][!A(x)])
    \cap \Lang L(\Gamma_2, \Delta_2)$ kacukupan.
    Miturut hipotesis induksi, kita nduwèni
    $\Lang L(!C_1) \subseteq \Lang L(\Gamma_1,
    \Delta_1, !A(c)) \cap \Lang L(\Gamma_2, \Delta_2)$.
    Dadi kudu dipriksa babagan~$c$: apa $c$ bisa ana ing~$!C_1$?
    Ora bisa. Yèn $c$ ana ing $!C_1$, banjur
    $c \in \Lang L(\Gamma_1, \Delta_1,
    !A(c))$ (iki bener) lan uga
    $c \in \Lang L(\Gamma_2, \Delta_2)$.
    Nanging iki ndadèkaké $c$ ana ing konklusi inferènsi
    $\RightR{\lforall}$ ing~$\pi$, nerak syarat variabel eigen.
    
    Kahanan kapindho padha polané.

    \item $\pi$ pungkasané nganggo $\RightR{\lexists}$
    kanthi !!{formula} utama~$\lexists[x][!A(x)]$.
    Ana rong kahanan manèh. Kita ngrembug nalika
    $\lexists[x][!A(x)]$ ana ing pérangan kapisan,
    dene kahanan kapindho dipasrahaké dadi latihan.

    Ing kahanan kapisan, $\pi$ pungkasané kaya mangkéné:
    \[
    \AxiomC{}
    \RightLabel{$\pi_1$}
    \Deduce$\maeh{\Gamma_1}{\Gamma_2} \fCenter
      \maeh{\Delta_1, \lexists[x][!A(x)], !A(t)}{\Delta_2}$
    \RightLabel{\RightR{\lexists}}
    \UnaryInf$\maeh{\Gamma_1}{\Gamma_2} \fCenter
      \maeh{\Delta_1, \lexists[x][!A(x)]}{\Delta_2}$
    \DisplayProof
    \]
    Hipotesis induksi maringi !!{proof}s kanggo
    $\Gamma_1 \Sequent \Delta_1, \lexists[x][!A(x)], !A(t), !C_1$
    lan $!C_1, \Gamma_2 \Sequent \Delta_2$.
    Kita bisa ngetrapaké $\RightR{\lexists}$ marang kang
    kapisan kanggo ngolèhké
    $\Gamma_1 \Sequent \Delta_1, \lexists[x][!A(x)], !C_1$.
    Mula $!C_1$ dadi interpolan yèn
    $\Lang L(!C_1) \subseteq \Lang L(\Gamma_1, \Delta_1,
    \lexists[x][!A(x)]) \cap \Lang L(\Gamma_2, \Delta_2)$
    kacukupan. Ing kéné hipotesis induksi mung maringi
    $\Lang L(!C_1) \subseteq
    \Lang L(\Gamma_1, \Delta_1, \lexists[x][!A(x)], !A(t))
    \cap \Lang L(\Gamma_2, \Delta_2)$. Élinga yèn kita
    nganggep ora ana !!{function}s. Mula suku~$t$ mung
    bisa !!a{constant}, yaiku $t \ident b$.
    Yèn $b \notin \Lang L(!C_1)$ utawa
    $b \in \Lang L(\Gamma_1, \Delta_1, \lexists[x][!A(x)]) \cap
    \Lang L(\Gamma_2, \Delta_2)$, $!C_1$ bisa langsung
    dijupuk dadi interpolan. Nanging kepriyé yèn
    $b \in \Lang L(!C_1)$ lan
    $b \notin \Lang L(\Gamma_1, \Delta_1,
    \lexists[x][!A(x)]) \cap \Lang L(\Gamma_2, \Delta_2)$?
    Sepisan, kanthi ngganti saben kedadéan~$b$ ing~$!C_1$
    nganggo !!a{variable}~$y$, kita bisa nulis
    $!C_1 \ident \Subst{!C_1'}{b}{y}$ kanthi $!C_1'$~ora
    ngemot~$b$. Sabanjuré, $b \notin
    \Lang L(\Gamma_1, \Delta_1, \lexists[x][!A(x)])$ utawa
    $b \notin \Lang L(\Gamma_2, \Delta_2)$.
    Kita bisa njupuk $!C \ident \lforall[y][!C_1'(y)]$
    utawa $!C \ident \lexists[y][!C_1'(y)]$ miturut urutané.
    Ing kahanan kapisan, kita nduwèni:
    \[
    \AxiomC{}
    \Deduce$\Gamma_1 \fCenter
      \Delta_1, \lexists[x][!A(x)], !A(b), !C_1'(b)$
    \RightLabel{\RightR{\lexists}}
    \UnaryInf$\Gamma_1 \fCenter \Delta_1, \lexists[x][!A(x)], !C_1'(b)$
    \RightLabel{\RightR{\lforall}}
    \UnaryInf$\Gamma_1 \fCenter
      \Delta_1, \lexists[x][!A(x)], \lforall[y][!C_1'(y)]$
    \DisplayProof
    \]
    lan
    \[
    \AxiomC{}
    \Deduce$!C_1'(b), \lforall[y][!C_1'(y)], \Gamma_2 \fCenter \Delta_2$
    \RightLabel{\LeftR{\lforall}}
    \UnaryInf$\lforall[y][!C_1'(y)], \Gamma_2 \fCenter \Delta_2$
    \DisplayProof
    \]
    !!{proof}s kanggo sekèn paling dhuwur asalé saka
    hipotesis induksi (ing kang kapindho, kita ngetrapaké
    admisibilitas pelemahan kanthi nambah
    $\lforall[y][!C_1'(y)]$ ing antesèden). Syarat
    variabel eigen kanggo \RightR{\lforall} ing
    !!{proof} kapisan kacukupan amarga
    $b \notin \Lang L(\Gamma_1, \Delta_1,
    \lexists[x][!A(x)])$, lan $!C_1'$~wis ditegesaké
    supaya ora ngemot~$b$.
    
    Ing kahanan kapindho, kita nduwèni (saka hipotesis
    induksi lan admisibilitas pelemahan kanthi nambah
    ~$\lexists[y][!C_1'(y)]$):
    \[
    \AxiomC{}
    \Deduce$\Gamma_1 \fCenter
      \Delta_1, \lexists[x][!A(x)], !A(b), \lexists[y][!C_1'(y)], !C_1'(b)$
    \RightLabel{\RightR{\lexists}}
    \UnaryInf$\Gamma_1 \fCenter
      \Delta_1, \lexists[x][!A(x)], \lexists[y][!C_1'(y)], !C_1'(b)$
    \RightLabel{\RightR{\lexists}}
    \UnaryInf$\Gamma_1 \fCenter
      \Delta_1, \lexists[x][!A(x)], \lexists[y][!C_1'(y)]$
    \DisplayProof
    \]
    lan
    \[
    \AxiomC{}
    \Deduce$!C_1'(b), \Gamma_2 \fCenter \Delta_2$
    \RightLabel{\LeftR{\lexists}}
    \UnaryInf$\lexists[y][!C_1'(y)], \Gamma_2 \fCenter \Delta_2$
    \DisplayProof
    \]
    Syarat variabel eigen kanggo \LeftR{\lexists} ing
    !!{proof} kapindho kacukupan amarga
    $b \notin \Lang L(\Gamma_2, \Delta_2)$, lan $!C_1'$
    wis ditegesaké supaya ora ngemot~$b$.
  \end{enumerate}
  Kahanan liyané padha polané; dipasrahaké dadi latihan.
\end{proof}

\begin{prob}
Anggepen ana konektif $\lnand$ kanthi aturan
\[
\Axiom$\Gamma \fCenter \Delta, !A$
\Axiom$\Gamma \fCenter \Delta, !B$
\RightLabel{\LeftR{\lnand}}
\BinaryInf$!A \lnand !B, \Gamma \fCenter \Delta$
\DisplayProof
\quad
\Axiom$!A, !B, \Gamma \fCenter \Delta$
\RightLabel{\RightR{\lnand}}
\UnaryInf$\Gamma \fCenter \Delta, !A \lnand !B$
\DisplayProof
\]
Tuduhna yèn Lema Maehara tetep bener kanthi ngrampungaké
kahanan-kahanan ing bukti \olref{lem:maehara} nalika
$\pi$ pungkasané nganggo salah siji aturan iki.
\end{prob}

Sawisé Lema Maehara kabuktèkaké, kita bisa migunakaké
kanggo mbuktèkaké Teorema Interpolasi Craig.

\begin{proof}[Bukti \olref{thm:interpolation}]
  Anggepen $!A \Sequent !B$ iku !!{provable}.
  Trapna \olref{lem:maehara} marang
  $\maeh{!A}{\ } \Sequent \maeh{\ }{!B}$ kanggo
  ngolèhké interpolan~$!C$. Kita nduwèni
  $\Proves !A \Sequent !C$ lan $\Proves !C \Sequent !B$.
\end{proof}

Anggepen $\Gamma$ iku teori (himpunan !!{sentence}s) ing
!!a{language}~$\Lang L$ kang ngemot predikat $n$-panggonan~$R$.
Kita kandha yèn $\Gamma$ \emph{nemtokaké kanthi implisit}~$R$
yèn lan mung yèn
\begin{align*}
\Gamma, \Gamma'
& \Entails \lforall[x_1][\dots\lforall[x_n][(R(x_1,\dots,x_n) \liff
R'(x_1, \dots, x_n))]],
\intertext{ing ngendi $\Gamma'$ asalé saka $\Gamma$ kanthi saben
kedadéan~$R$ diganti~$R' \notin \Lang L$. $\Gamma$
\emph{nemtokaké kanthi eksplisit}~$R$ yèn lan mung yèn ana~$!A(x_1, \dots, x_n)$
kang ora ngemot~$R$ saéngga}
\Gamma &\Entails
\lforall[x_1][\dots\lforall[x_n][(R(x_1,\dots,x_n) \liff !A(x_1,
\dots, x_n))]].
\end{align*}
Cetha yèn $\Gamma$ nemtokaké $R$ kanthi eksplisit,
mula uga nemtokaké kanthi implisit~$R$.
Waliké ora langsung cetha, nanging tetep bener:

\begin{lem}[Teorema Definabilitas Beth]
  Yèn $\Gamma$ nemtokaké~$R$ kanthi implisit,
  teori iku uga nemtokaké kanthi eksplisit.
\end{lem}

\begin{proof}
  Anggepen $\Gamma$ nemtokaké~$R$ kanthi implisit.
  Upama $R'$ lan $\Gamma'$ kaya ing ndhuwur, sarta
  $\Lang L' = \Lang L(\Gamma') = \Lang L \setminus \{R\}
  \cup \{R'\}$. Miturut hipotesis, kita nduwèni
  \begin{align*}
    \Gamma, \Gamma' & 
    \Sequent \lforall[x_1][\dots\lforall[x_n][(R(x_1,\dots,x_n) 
    \liff R'(x_1, \dots, x_n))]]
  \intertext{Upama $c_1$, \dots,~$c_n$ iku !!{constant}s kang ora ana ing~$\Gamma$. Miturut lema inversibilitas, kita nduwèni}
    \Gamma, \Gamma', R(c_1,\dots,c_n) & \Sequent  R'(c_1, \dots, c_n)
  \intertext{Trapna Lema Maehara marang}
    \maeh{\Gamma, R(c_1,\dots,c_n)}{\Gamma'} & \Sequent
      \maeh{\ }{R'(c_1, \dots, c_n)}
  \intertext{kanggo nemokaké~$!A \ident !A(x_1, \dots, x_n)$ saéngga}
    \Gamma, R(c_1,\dots,c_n) & \Sequent !A(c_1, \dots, c_n) \\
    !A(c_1, \dots, c_n), \Gamma' & \Sequent R'(c_1, \dots, c_n)
  \intertext{nduwèni !!{proof}s. Amarga $\Lang L(!A) \subseteq \Lang L_1 \cap
  \Lang L_2$, $!A$~ora ngemot~$R$, kang mung ana ing pérangan kapisan.
  Gantenana $R'$ nganggo~$R$ ing saindhenging !!{proof} sekèn kapindho
  kanggo ngolèhké !!a{proof} kanggo}
    !A(c_1, \dots, c_n), \Gamma & \Sequent R(c_1, \dots, c_n)
  \end{align*}
  Trapna \RightR{\lif} lan \RightR{\lforall} kanggo ngolèhké
  !!a{proof} kanggo
  \[
    \Gamma \Sequent \lforall[x_1][\dots\lforall[x_n][(R(x_1,\dots,x_n) \liff !A(x_1, \dots, x_n))]].
  \]
  Iki nuduhaké yèn $!A(x_1, \dots, x_n)$ nemtokaké
  $R$ kanthi eksplisit ing~$\Gamma$.
\end{proof}

Asil katelu kang équivalèn karo interpolasi (lan kaya
interpolasi, bisa dibuktèkaké nganggo Lema Maehara)
yaiku teorema iki: Yèn rong teori konsistèn~$\Gamma_1$
lan $\Gamma_2$ ngemot teori komplet $\Gamma$ ing basa
$\Lang L(\Gamma) \subseteq \Lang L(\Gamma_1)
\cap \Lang L(\Gamma_2)$, banjur
$\Gamma_1 \cup \Gamma_2$ konsistèn.

Élinga yèn teori komplet iku teori kang kanggo saben
!!{sentence}~$!A$ ing basané, $\Gamma \Proves !A$ utawa
$\Gamma \Proves \lnot !A$. Amarga $\Gamma_1$ lan~$\Gamma_2$
iku ekstènsi konsistèn saka~$\Gamma$, mula $\Gamma$
uga konsistèn. Dadi, yèn $!A$ iku !!a{sentence} ing basa
$\Lang L(\Gamma_1) \cap \Lang L(\Gamma_2)$, ora mungkin
$\Gamma_1 \Entails !A$ lan
$\Gamma_2 \Entails \lnot !A$ bebarengan.
Anggepen ana~$!A$ mangkono. Yèn
$\Gamma_1 \Entails !A$, banjur $\Gamma \Entails !A$:
yèn ora, amarga $\Gamma$ komplet,
$\Gamma \Entails \lnot !A$, lan amarga
$\Gamma \subseteq \Gamma_1$, kita ngolèhké
$\Gamma_1 \Entails \lnot !A$, kang ndadèkaké
$\Gamma_1$ ora konsistèn. Dadi $\Gamma \Entails !A$,
lan amarga $\Gamma \subseteq \Gamma_2$,
$\Gamma_2 \Entails !A$. Mula $\Gamma \Entails/
\lnot !A$; yèn ora, $\Gamma_2$ ora konsistèn.

Ora anané !!a{sentence}~$!A$ ing
!!{language} sesarengan~$\Lang L(\Gamma_1) \cap \Lang L(\Gamma_2)$
kanthi $\Gamma_1 \Entails !A$ lan
$\Gamma_2 \Entails \lnot !A$ wis cukup kanggo bukti,
kang ing kéné kita wènèhaké kanthi cara teori bukti
nganggo Lema Maehara.

\begin{thm}[Teorema Konsistènsi Gabungan Robinson]
Anggepen $\Gamma_1$ lan $\Gamma_2$ iku teori konsistèn,
lan ora ana~$!A$ ing !!{language}
$\Lang L(\Gamma_1) \cap \Lang L(\Gamma_2)$ kang marakaké
$\Gamma_1 \Entails !A$ lan
$\Gamma_2 \Entails \lnot !A$ bebarengan. Banjur
$\Gamma_1 \cup \Gamma_2$ konsistèn.
\end{thm}

\begin{proof}
  Kosok balènana, anggepen $\Gamma_1 \cup \Gamma_2$
  ora konsistèn. Banjur sekèn
  $\Gamma_1, \Gamma_2 \Sequent \ $ nduwèni !!a{proof}.
  Trapna Lema Maehara marang
  $\maeh{\Gamma_1}{\Gamma_2} \Sequent \maeh{\ }{\ }$
  kanggo ngolèhké !!a{sentence}~$!A$ ing
  !!{language}~$\Lang L(\Gamma_1) \cap \Lang L(\Gamma_2)$
  saéngga $\Proves \Gamma_1 \Sequent !A$ lan
  $\Proves !A, \Gamma_2 \Sequent \quad$. Saka sekèn
  kapindho iki kita ngolèhké
  $\Proves \Gamma_2 \Sequent \lnot !A$.
  Nanging kita wis nganggep ora ana !!{sentence} mangkono.
\end{proof}

Ing ndhuwur kita meneng-menengan nganggep teori
$\Gamma$, $\Gamma_1$, lan $\Gamma_2$ winates.
Miturut kakompakan, asilé uga bener kanggo teori
tanpa wates.

\end{document}
